\clearpage\fancyhead[L]{\small PALLAS. COMET.}\hypertarget{native-332-lower}{}\label{native-332-lower}
\begin{center}\textbf{PALLAS. COMET.}\par\vspace{4pt}\hrule\vspace{4pt}\fontsize{10}{12}\selectfont Göttingische gelehrte Anzeigen. Stück 130. Seite 1289\,..\,1293. 1811. August 17.\par\vspace{4pt}\hrule\vspace{4pt}\end{center}

\par\indent Im 8. Stück unserer Anzeigen von d. J. [S. 322 d. B.] ist von der Arbeit des Hrn. Prof. \textsc{Gauss} über die Störungen der Pallas seit ihrer ersten Entdeckung Nachricht gegeben. Auch ist daselbst die \SourceLeaf{333}erste, am 4. Januar auf der hiesigen Sternwarte gemachte Beobachtung mitgetheilt, wodurch die Resultate jener Arbeit mit einer kaum erwarteten Genauigkeit bestätigt wurden. Da der Planet zu lichtschwach war, um auf unserer Sternwarte im Meridian beobachtet werden zu können, so wurden die Beobachtungen hier nicht weiter fortgesetzt, sondern nur auswärtige Astronomen, denen stärkere Instrumente zu Gebote stehen, zeitig genug benachrichtigt. So hat Hr. Prof. \textsc{Gauss} für die Zeit der Opposition, welche im Februar eintrat, von fünf Sternwarten her, nemlich von Paris, Mannheim, Seeberg, Berlin und Hamburg, sehr schätzbare Beobachtungen erhalten, die ihm zur Grundlage seiner weitern Untersuchungen gedient haben. Die auf der Seebergen und Mannheimer Sternwarte von den Herren \textsc{von Lindenau}, \textsc{Barry} und \textsc{Harding} angestellten Beobachtungen sind bereits in der Monatlichen Correspondenz abgedruckt, die übrigen theilen wir hier mit.

\Needspace{14\baselineskip}\begin{center}\fontsize{10}{12}\selectfont\setlength{\tabcolsep}{4pt}\renewcommand{\arraystretch}{1.15}
Beobachtungen der Pallas zu Paris. von Hrn. \textsc{Bouvard}.

\begin{tabular}{c|c|c}
im Meridian & Gerade Aufsteig. & Abweichung südl. \\
\hline
1811 Febr. 17 & $146^\circ\,58'\,33''15$ & $13^\circ\,29'\,42''4$ \\
18 & $146\;\;47\;\;55.95$ & $13\;\;\;5\;\;59.4$ \\
22 & $146\;\;\;6\;\;27.90$ & $11\;\;28\;\;19.6$ \\
27 & $145\;\;18\;\;13.65$ & $9\;\;19\;\;32.9$ \\
\end{tabular}
\end{center}

\Needspace{14\baselineskip}\begin{center}\fontsize{10}{12}\selectfont\setlength{\tabcolsep}{4pt}\renewcommand{\arraystretch}{1.15}
Beobachtungen zu Berlin, von Hrn. Prof. \textsc{Bode}.

\begin{tabular}{c|c|c|c}
1811 & Mittlere Zeit & Gerade Aufst. & Abweichung südl. \\
\hline
Febr. 18 & $11^{\mathrm u}\,55^{\mathrm m}\,22^{\mathrm s}5$ & $146^\circ\,48'\,2''$ & --- --- --- \\
19 & $11\;\;50\;\;34.0$ & $146\;\;37\;\;32$ & $12^\circ\,42'\,44''5$ \\
20 & $11\;\;45\;\;57.0$ & $146\;\;27\;\;\;4$ & $12\;\;19\;\;\;1$ \\
März 12 & $10\;\;16\;\;33.6$ & $143\;\;45\;\;29$ & $3\;\;32\;\;42$ \\
16 & $9\;\;59\;\;45.1$ & $143\;\;29\;\;18$ & $1\;\;48\;\;\;3$ \\
17 & $9\;\;55\;\;35.1$ & $143\;\;26\;\;13$ & $1\;\;22\;\;11$ \\
\end{tabular}
\end{center}

\Needspace{14\baselineskip}\begin{center}\fontsize{10}{12}\selectfont\setlength{\tabcolsep}{4pt}\renewcommand{\arraystretch}{1.15}
Beobachtungen zu Hamburg, von Hrn. Professor \textsc{Schumacher}.

\begin{tabular}{c|c|c}
im Meridian & Gerade Aufsteig. & Abweichung südl. \\
\hline
1811 Febr. 19 & $146^\circ\,37'\,15''0$ & $12^\circ\,42'\,45''3$ \\
20 & $146\;\;26\;\;52.5$ & $12\;\;18\;\;32.7$ \\
21 & $146\;\;16\;\;37.5$ & $11\;\;53\;\;47.7$ \\
22 & $146\;\;\;6\;\;36.0$ & $11\;\;28\;\;53.5$ \\
März 9 & $144\;\;\;1\;\;45.0$ & $4\;\;52\;\;42.5$ \\
15 & $143\;\;32\;\;39.0$ & $2\;\;13\;\;59.0$ \\
16 & $143\;\;29\;\;24.0$ & $1\;\;47\;\;57.0$ \\
18 & $143\;\;23\;\;34.5$ & $0\;\;56\;\;47.8$ \\
\end{tabular}
\end{center}


\Needspace{8\baselineskip}\noindent Hieraus leitete Hr. Prof. \textsc{Gauss} folgendes Resultat für die Opposition ab:
\begin{center}1811 Febr. 21. $19^{\mathrm u}\,25^{\mathrm m}\,31^{\mathrm s}$ mittlere Zeit in Göttingen\end{center}

\begin{center}\fontsize{10}{12}\selectfont\setlength{\tabcolsep}{4pt}\renewcommand{\arraystretch}{1.15}\begin{tabular}{llr}
Wahre heliocentrische Länge & .\;.\;.\;.\;.\;.\;.\;. & $152^\circ\,48'\,15''8$\\
Wahre geocentrische Breite, südlich & .\;.\;.\;.\;. & $23\;48\;19.2$\\
\end{tabular}\end{center}

\Needspace{17\baselineskip}\par\indent Die Übereinstimmung dieser Opposition mit den sechs vorhergehenden, nach gehöriger Anbringung der Störungen von dem Jupiter, wie Hr. Prof. \textsc{Gauss} sie entwickelt hatte, ist in der That bewundernswürdig. Nachdem nur einige äusserst unbedeutende Correctionen an die schon vor der Auffindung bestimmten Elemente angebracht waren, ergaben sich die Unterschiede zwischen der Beobachtung und Rechnung, wie folgt:
\SourceLeaf{334}
\begin{center}\fontsize{10}{12}\selectfont\setlength{\tabcolsep}{4pt}\renewcommand{\arraystretch}{1.15}
Unterschied der

\begin{tabular}{c|c|c}
 & mittleren Länge & heliocentrischen Breite \\
\hline
1803 & $+8''7$ & $-1''1$ \\
1804 & $-7.0$ & $-10.2$ \\
1805 & $+1.4$ & $-3.5$ \\
1807 & $-1.2$ & $-2.4$ \\
1808 & $-3.1$ & $-21.4$ \\
1809 & $+4.5$ & $-2.1$ \\
1811 & $-2.2$ & $-12.1$ \\
\end{tabular}
\end{center}


\Needspace{21\baselineskip}\par\indent Hiernächst wurde die Berechnung der Störungen noch ein Jahr weiter fortgesetzt. Hr. Prof. \textsc{Gauss} wurde bei dieser beschwerlichen Arbeit unterstützt von Hrn. \textsc{Nicolai}, welcher sich bei uns den mathematischen Wissenschaften mit grossem Eifer widmet, und unter Hrn. Prof. \textsc{Gauss} Aufsicht den grössern Theil der numerischen Rechnungen mit eben so viel Fleiss als Geschicklichkeit ausgeführt hat. Den Resultaten dieser Rechnung zufolge, wird die Bewegung der Pallas um die Zeit der Opposition des nächsten Jahres sich durch folgende elliptische, die Störungen bereits einschliessende, Elemente darstellen lassen:

\begin{center}\fontsize{10}{12}\selectfont\setlength{\tabcolsep}{4pt}\renewcommand{\arraystretch}{1.15}\begin{tabular}{llr}
\shortstack[l]{Epoche der mittlern Länge 1812 Junius 10 Mittag\\im Meridian von Göttingen} & . & $239^\circ\,4'\,46''1$\\
Länge des Perihels für dieselbe Zeit &  & $121\;0\;48.5$\\
Länge des aufsteigenden Knoten & . & $172\;32\;44.5$\\
Neigung der Bahn & .\;.\;.\;.\;.\;.\;. & $34^\circ\,34'\,54''7$\\
Tägliche mittlere tropische Bewegung &  & $768''5746$\\
Excentricität & .\;.\;.\;.\;.\;. & $=\sin13^\circ\,59'\,1''8$\\
Logarithm. der halben grossen Axe & . & $0.4429321$\\
\end{tabular}\end{center}

\Needspace{20\baselineskip}\par\indent Nach eben diesen Elementen hat Hr. \textsc{Nicolai} auch eine Ephemeride für die Erscheinung der Pallas im Jahre 1812 berechnet, die ihrer Ausdehnung wegen hier nicht Platz finden kann, und in der Monatl. Correspondenz und im Astronom. Jahrbuche abgedruckt werden wird. Wir setzen nur noch die gleichfalls von Hrn. \textsc{Nicolai} \emph{im voraus} berechnete nächste Opposition (die achte) her, die zu ihrer Zeit zur Prüfung der Genauigkeit unserer Resultate wird dienen können:
\begin{center}1812 Juni 10. $3^{\mathrm u}\,32^{\mathrm m}$ mittlere Zeit in Göttingen.\end{center}

\begin{center}\fontsize{10}{12}\selectfont\setlength{\tabcolsep}{4pt}\renewcommand{\arraystretch}{1.15}\begin{tabular}{llr}
Heliocentrische Länge der Pallas & .\;.\;.\;.\;.\;.\;.\;. & $259^\circ\,28'\,53''$\\
Heliocentrische Breite, nordlich & .\;.\;.\;.\;.\;.\;.\;. & $34\;32\;37$\\
Geocentrische Breite & .\;.\;.\;.\;.\;.\;.\;.\;.\;.\;.\;. & $48\;16\;36$\\
\end{tabular}\end{center}

\begin{center}\fontsize{10}{12}\selectfont\setlength{\tabcolsep}{4pt}\renewcommand{\arraystretch}{1.15}\begin{tabular}{lr}
\shortstack[l]{Logarithm. der Entfernung der Pallas\\von der Sonne} & $0.504259$\\
\shortstack[l]{Logarithm. der Entfernung der Pallas\\von der Erde} & $0.384915$\\
\end{tabular}\end{center}

\Needspace{20\baselineskip}\par\indent Bei dieser Gelegenheit fügen wir noch die parobolischen  Elemente des im März dieses Jahres von \textsc{Flaugergues} entdeckten Cometen bei, welche Hr. Prof. \textsc{Gauss} soeben aus einigen Beobachtungen des Hrn. v. \textsc{Zach} zu Marseille vorläufig berechnet hat, und die zur Erleichterung der Wiederauffindung dieses merkwürdigen Cometen gegen Ende August werden dienen können:
\begin{center}Durchgang durchs Perihelium 1811 Sept. 10. $0^{\mathrm u}\,51^{\mathrm m}$\end{center}

\begin{center}\fontsize{10}{12}\selectfont\setlength{\tabcolsep}{4pt}\renewcommand{\arraystretch}{1.15}\begin{tabular}{llr}
Länge des Perihelium & .\;.\;.\;.\;.\;.\;.\;.\;.\;.\;.\;. & $73^\circ\,14'\,35''$\\
Logarithm. des kleinsten Abstandes & .\;.\;.\;.\;.\;.\;. & $9.99153$\\
Länge des aufsteigenden Knoten & .\;.\;.\;.\;.\;.\;.\;. & $141^\circ\,4'\,59''$\\
Neigung der Bahn & .\;.\;.\;.\;.\;.\;.\;.\;.\;.\;.\;.\;. & $73\;48\;2$\\
Bewegung \emph{rückläufig} &  & \\
\end{tabular}\end{center}

\SourceLeaf{335}\Needspace{12\baselineskip}\par\indent Folgende Stellungen sind nach diesen Elementen berechnet:
\begin{center}\fontsize{10}{12}\selectfont\setlength{\tabcolsep}{4pt}\renewcommand{\arraystretch}{1.15}
\begin{tabular}{l|cc|c}
Mitternacht & Gerade Aufst. & Nordl. Abw. & Lichtstärke \\
\hline
1811 Aug. 11 & $142^\circ\,43'$ & $29^\circ\,54'$ & $0.194$ \\
\phantom{1811 }\phantom{Aug. }26 & $152\;\;19$ & $36\;\;15$ & $0.294$ \\
Sept. 10 & $167\;\;34$ & $43\;\;14$ & $0.434$ \\
\phantom{Sept. }25 & $194\;\;22$ & $48\;\;35$ & $0.576$ \\
Octob. 10 & $232\;\;28$ & $44\;\;52$ & $0.595$ \\
\end{tabular}
\end{center}


\par\indent Nach demselben Massstabe war die Lichtstärke (unter der Voraussetzung, dass der Comet sein Licht blos von der Sonne erhält) am 11. April $0.035$, und am 2. Juni $0.052$; er wird also zu Anfang Octobers 17 mal so hell sein, als am 11. April, und sich daher in ganz vorzüglichem Glanze am nordlichen Himmel zeigen. Nachher wird er noch mit abnehmendem Lichte die Sternbilder des Herkules und des Adlers durchlaufen, und im Anfange des nächsten Jahres im Wassermann unsichtbar werden. Hindert sein dann allmählig gar zu schwach gewordenes Licht es nicht, so kann er vielleicht im April noch einmal im Wassermann wieder sichtbar, und mit sehr starken Telescopen noch eine geraume Zeit verfolgt werden. Aber wenn dies auch nicht gelingen sollte, so wird er doch wegen seines ausgezeichneten Glanzes im Herbste dieses Jahrs und wegen seiner langen Sichtbarkeit, die gewiss an zehn Monate betragen wird, in den Annalen der Cometographie eine der allermerkwürdigsten Erscheinungen bleiben.
